% Source: Weinberg GR, physical PDF pp. 273--274 (printed pp. 251--252).
% Coverage: Chapter 10 opening material before Section 10.1.
% Figures/tables/footnotes: chapter epigraph; no figures or tables.
% Status: source-reviewed and compile-clean.
% Uncertainties: none.

\begin{flushright}
\begin{minipage}{0.42\textwidth}
\itshape
``It is the nature of all\\
greatness not to be exact.''\\
Edmund Burke, speech on\\
American taxation, 1774
\end{minipage}
\end{flushright}

We have seen a great many similarities between gravitation and
electromagnetism. It should therefore come as no surprise that Einstein's
equations, like Maxwell's equations, have radiative solutions.

No one has yet certainly detected gravitational radiation, but the reason for
this is not hard to find; Einstein's theory predicts that gravitational
radiation is produced in extremely small quantities in ordinary atomic
processes. For instance, the probability that a transition between two atomic
states will proceed by emission of gravitational, rather than electromagnetic,
radiation is typically of order \(GE^2/e^2\), where \(E\) is the energy
released and \(e\) is the electronic charge. For \(E\simeq1\ \mathrm{eV}\),
this probability is about \(3\times10^{-54}\).

Why then study gravitational radiation? One reason is of course that some day
we may find a strong source of gravitational radiation. Such a source may
indeed already have been detected. (See Section~\ref{sec:10.7}.) However,
gravitational radiation would be interesting even if there were no chance of
ever detecting any, for the theory of gravitational radiation provides a
crucial link between general relativity and the microscopic frontier of
physics.

We have learned in recent years to describe the fundamental observables of
microscopic phenomena in terms of elementary particles and their collisions.
In classical electrodynamics it is the plane-wave solutions of Maxwell's
equations that lead most naturally to an interpretation in terms of a
particle, the photon. Similarly, it is the radiative solutions of Einstein's
equations that will lead here to the concept of a particle of gravitational
radiation, the \emph{graviton}.

Unfortunately, the theory of gravitational radiation is complicated by the
nonlinearity of Einstein's equations. In the spirit of Section~\ref{sec:7.6},
we may say that any gravitational wave is itself a distribution of energy and
momentum that contributes to the gravitational field of the wave. This
complication prevents our being able to find general radiative solutions of
the exact Einstein equations.

There are two approaches to this difficulty. One is to study only the
weak-field radiative solutions of Einstein equations, which describe waves
carrying not enough energy and momentum to affect their own propagation. The
other approach is to look long and hard for special solutions of the exact
Einstein field equations. A great deal of mathematical ingenuity has gone
into the second approach with results of some elegance. However, this chapter
deals with only the first, weak-field, approach to gravitational radiation.
One reason is that any observable gravitational radiation is likely to be of
very low intensity. A second deeper reason is that it is only possible to
attach a precise meaning to the concept of an elementary particle when it is
far away from all other particles, and for gravitons this corresponds to a
weak-field solution of the field equations.

The reader should not conclude that there is any fundamental gap in our
understanding of gravitation because we cannot find general exact solutions of
the nonlinear field equations. Indeed, similar problems arise in
electrodynamics. The problem of computing the exact electromagnetic field
produced by a decaying current in an electrical oscillator is highly
nonlinear, because the field acts back on the current that produces it. Even
though this problem was not solved for many years after Maxwell's theory,
still there was no doubt that electrical oscillators would produce the
electromagnetic waves studied by Maxwell. Gravitational waves are more
complicated than electromagnetic waves because they contribute to their own
source outside the material gravitational antenna. However, the simple
properties of both electromagnetic and gravitational waves emerge when we
look far out into the wave zone, where the fields are weak.
